\documentclass{article}
\usepackage{amsmath}
\usepackage{amssymb}
\usepackage[margin=1in]{geometry}

\begin{document}

\section*{Math Notes: Residue Classes mod 6}

\subsection*{Residue classes}

Every integer has one residue modulo 6:
\[
n \equiv 0,1,2,3,4,5 \pmod{6}.
\]

\subsection*{Prime residue constraint}

For every prime greater than 3:
\[
p > 3 \Rightarrow p \equiv \pm 1 \pmod{6}.
\]

Equivalently:
\[
p > 3 \Rightarrow p \equiv 1 \text{ or } 5 \pmod{6}.
\]

\subsection*{Necessary but not sufficient}

The residue constraint is necessary:
\[
p>3 \Rightarrow p \equiv \pm 1 \pmod{6}.
\]

It is not sufficient:
\[
n \equiv \pm 1 \pmod{6} \nRightarrow n \text{ is prime}.
\]

\subsection*{CGCS score}

For this notebook:
\[
CGCS_{mod6} =
\frac{\#\{p>3 : p \equiv 1,5 \pmod{6}\}}
{\#\{p>3\}}.
\]

Measured result:
\[
CGCS_{mod6} = 1.000000.
\]

\subsection*{Drift}

\[
drift_{mod6} = 1 - CGCS_{mod6}.
\]

Measured result:
\[
drift_{mod6} = 0.000000.
\]

\subsection*{Recoverability}

Modulo 6 recovers candidate residue classes:
\[
n \equiv \pm 1 \pmod{6}.
\]

It does not recover prime identity exactly.

\end{document}
